\section{Additional details about the \texttt{Discover-and-Cover} algorithm}\label{sec:app_algo}


In this section, we provide all the details about the \texttt{Discover-and-Cover} (Algorithm~\ref{alg:final_algorithm}) algorithm that are omitted from the main body of the paper.
%
In particular:
\begin{itemize}
	\item Appendix~\ref{sec:app_params} gives a summary of the definitions of the parameters required by Algorithm~\ref{alg:final_algorithm}, which are set as needed in the proofs provided in the rest of this section.
	\item Appendix~\ref{sec:app_action_oracle} provides the proofs of the lemmas related to the \texttt{Action-Oracle} procedure.
	\item Appendix~\ref{sec:app_try_partiton} provides the proofs of the lemmas related to the \texttt{Try-Cover} procedure.
	\item Appendix~\ref{sec:app_fiind_contr} provides the proofs of the lemmas related to the \texttt{Find-Contract} procedure.
	\item Appendix~\ref{sec:app_proof_thm2} provides the proof of the final result related to Algorithm~\ref{alg:final_algorithm} (Theorem~\ref{thm:finalthm}).
\end{itemize}

\allowdisplaybreaks

\subsection{Definitions of the parameters in Algorithm~\ref{alg:final_algorithm}}\label{sec:app_params}

The parameters required by Algorithm~\ref{alg:final_algorithm} are defined as follows:
\begin{itemize}
	\item $\displaystyle \epsilon  \coloneqq \frac{\rho^2}{32^2 B m^2 n^2}  $
	\item $\displaystyle \eta \coloneqq \frac{\epsilon \sqrt{m} n }{2}$
	\item $\displaystyle \alpha \coloneqq \frac{\delta}{2 n^3 \left[ \log \left(\frac{2Bm}{\eta}\right) + \binom{ m + n +1}{m}  \right]}$
	\item $\displaystyle q \coloneqq \left\lceil \frac{1}{2 \epsilon^2} \log \left(\frac{2 m}{\alpha}\right) \right\rceil$
\end{itemize}

\subsection{Proofs of the lemmas related to the \texttt{Action-Oracle} procedure}\label{sec:app_action_oracle}



	%As a first step, we show that given the meta-action $d$, it holds $\lVert\widetilde F_1-\widetilde F_2\rVert_{\infty}\le 4 \epsilon n$ for each $  \widetilde F_1,\widetilde F_2\in \mathcal{F}[d]$.
	%
	%
	\begin{comment}
	Let $I(d) \subseteq \mathcal{A}$ be the set of actions such that Algorithm~\ref{alg:action_oracle} returned the meta-action $d$ (for a contract $p'$).
%
As a first step, we show that \begin{align} \label{eq:rec}
	\lVert F_{a_i}- F_{a_j}\rVert_{\infty}\le 4 \epsilon n
\end{align}
for each $a_i$, $a_j \in I(d)$.
%
We prove this result with an induction on the set $I(d)$.
The size of $I(d)$ can increase in two ways: i) Algorithm~\ref{alg:action_oracle} returns $d$ when the best response is an action $\hat a \notin I(d)$; ii) Algorithm~\ref{alg:action_oracle} merges a set of meta-actions $\{d_i\}_{i \in \mathcal{I}}$.
%For the ease of exposition, we let $I(d)=\{a_1,\ldots,a_q\}$, where the index denote the order in which the actions are added to $I(d)$ (in a sequence of different calls to Algorithm~\ref{alg:action_oracle}).
Consider case i). We show that for action $\hat a$, there exists another action $\bar a \in I(d)$ such that $\lVert  F_{\hat a}- F_{\bar a} \rVert \le 4 \epsilon$.
Let $\widetilde F_{\hat a}$ be the distribution computed by Algorithm~\ref{alg:action_oracle} when $\hat a$ is added to $I(d)$. 
%
Since $\hat a_i$ is added to $I(d)$, there exists a distribution $\widetilde F_{\bar a}  \in \mathcal F[d]$ such that $\lVert\widetilde F_{\hat a}-\widetilde F_{\bar a} \rVert\le 2\epsilon$, where we let  $\bar a \in I(d)$ be the action played by the agent when the distribution $\widetilde F_{\bar a}$ is computed by Algorithm~\ref{alg:action_oracle}. 
Under the clean event $\mathcal{E}_\epsilon$, we have that $\lVert \widetilde F_{\hat a}- F_{\hat a}\rVert_\infty\le \epsilon$ and $\lVert \widetilde F_{\bar a}- F_{\bar a}\rVert_\infty\le \epsilon$. This implies that $\lVert F_{\hat a}- F_{\bar a} \rVert\le 4\epsilon$.

Consider case ii). We use an argument similar to the one in i).
Let $\widetilde F_{\hat a}$ be the distribution computed by Algorithm~\ref{alg:action_oracle} when the set of meta-actions $\{d_i\}_{i \in \mathcal{I}}$ are merged, where  $\hat a$ is the best response to the play contract. 
Since $\{d_i\}_{i \in \mathcal{I}}$ are merged, for each $d_i$, $i \in \mathcal{I}$, there exists a distribution $\widetilde F_{\bar a}  \in \mathcal F[d_i]$ such that $\lVert\widetilde F_{\hat a}-\widetilde F_{\bar a} \rVert\le 2\epsilon$, where we let  $\bar a \in I(d_i)$ be the action played by the agent when the distribution $\widetilde F_{\bar a}$ is computed by Algorithm~\ref{alg:action_oracle}. 
Under the clean event $\mathcal{E}_\epsilon$, we have that $\lVert \widetilde F_{\hat a}- F_{\hat a}\rVert_\infty\le \epsilon$ and $\lVert \widetilde F_{\bar a}- F_{\bar a}\rVert_\infty\le \epsilon$. This implies that $\lVert F_{\hat a}- F_{\bar a} \rVert\le 4\epsilon$.

Hence, by induction we have proved Equation~\eqref{eq:rec}.
Since there are at most $n$ actions in $I(d)$, it follows that for every two actions $a,a' \in I(d) $, it holds   
\[\lVert F_{a}- F_{a'} \rVert\le 4n\epsilon.\]

Then, since the clean event $\mathcal{E}_\epsilon$ holds,  we note that \[\lVert \widetilde F_d -F_{\hat a} \rVert_\infty\le \epsilon,\]
where $\hat a \in I(d)$ is the action played by the agent when the distribution $\widetilde F_d $ is computed by Algorithm~\ref{alg:action_oracle}. 
%

%Moreover, under the event $\mathcal{E}_\epsilon$, for every $\widetilde{F} \in \mathcal{F}[d]$ there always exists an action $a \in I(d)$ such that $\| \widetilde{F} - F_{a}\|_{\infty} \le \epsilon$. 

Let $a^*\coloneqq a^*(p)$ be the best response to contract $p$. 
Combining the two inequalities, we get that
\begin{equation*}
	\|\widetilde F_{d}- F_{a^*}\|_{\infty} \leq \|\widetilde F_d- F_{\hat a} \|_{\infty} + \|F_{\hat a}- F_{a^*} \|_{\infty}  \le 4\epsilon n+ \epsilon\le 5\epsilon n.
\end{equation*}
This implies that $a^*(p) \in A(d)$, concluding the proof.
	\end{comment}

%	Finally, let $f_d$; close to $F_a$, ends.
	
%	For each distribution $F \in \mathcal{F}[d]$
%	
%	Algorithm~\ref{alg:action_oracle} guarantees that since $  \tilde F_1,\tilde F_2\in \mathcal{F}[d]$, there exists a sequence of $\tilde F_1,\ldots,\tilde F_q$ of probability distributions such that $\lVert \tilde F_i- \tilde F_{i-1} \rVert_{\infty}\le 2\epsilon$ for each $i \in \{2 \ldots,q\}$, where we let $\tilde F_q \coloneqq \tilde F_2$.
%	
%	
%	%
%	For ease of presentation, in the rest of this proof we assume that agent's actions are indexed by natural numbers, and denote by $a_i, a_j$ two different actions in $\mathcal{A}$.
%	% 
%	Moreover, in the following we fix a meta-action $d \in \mathcal{D}$, as the proof will hold for any choice of $d$.
%	
%	In order to prove the lemma, we introduce the following iterative procedure.
%	%
%	Let $a_i \in \mathcal{A}$ be the agent's action associated to the empirical distribution $\widetilde{F}_d$ we initialize a set $\mathcal{J}\gets \{a_i\}$.
%	%
%	Then, proceed by iteratively adding to $\mathcal{J}$ all the actions $a_i$ in $B(d) \setminus \mathcal{J}$ that satisfy the condition $\| F_{a_i} - F_{a_j} \|_{\infty} \leq 4 \epsilon$ for at least one $a_j$ already in $\mathcal{J}$, until there are no more such actions.
%	%
%	We show that the iterative procedure described above eventually terminates with $\mathcal{J}$ equal to $B(d)$, and this allows us to easily prove the statement of the lemma.
%	%
%	% Then we show that through this process, we will eventually have $\mathcal{J}$ equal to $I(d)$.
%	
%	Intuitively, the procedure outlined above ensures that each action $a_i$ added to $\mathcal{J}$ induces a distribution over outcomes $F_{a_i}$ whose distance with respect to another distribution $F_{a_j}$ induced by an action $a_j $ already in $\mathcal{J}$ is bounded.
%	%
%	This guarantees that the set $\mathcal{J}$ contains actions whose distributions over outcomes are ``sufficiently close'' to each other.
%	%
%	Then, by showing that $\mathcal{J}$ eventually coincides with the set $B(d)$, we can guarantee that such a ``closeness'' property holds for the whole set $B(d)$. 
%	
%	
%	As a crucial step in proving the desired property of the procedure described above, we show that, at any point of the procedure, either there exists an action $a_i \in I(d) \setminus \mathcal{J}$ such that $\| F_{a_i} - F_{a_j}\|_{\infty} \le 4 \epsilon $ for some $a_j \in \mathcal{J}$ or $ B(d) \setminus \mathcal{J}$ is an empty set.
%	%
%	In the case in which $ I(d) \setminus \mathcal{J}$ in non-empty, given how Algorithm~\ref{alg:action_oracle} works, under the event $\mathcal{E}_\epsilon$ there must exist two empirical distributions $\widetilde{F},\widetilde{F}' \in \mathcal{F}[d]$ such that: (i) $\|\widetilde{F} - F_{a_j}\|_{\infty} \leq \epsilon$ for some $a_j \in \mathcal{J}$, (ii) $\|\widetilde{F}' - F_{a_i} \|_{\infty} \leq \epsilon$ for some $a_i \in I(d)\setminus \mathcal{J}$, and (iii) $\|\widetilde{F}-\widetilde{F}'\|_{\infty}  \leq 2\epsilon$. This is because all the actions belonging to either $\mathcal{J}$ or $I(d) \setminus \mathcal{J}$ are identified as a single meta-action $d$ by the algorithm.
%	%
%	As a consequence, either $I(d) \setminus \mathcal{J}$ is an empty set or, by employing the triangular inequality, there exists an action $a_i \in I(d) \setminus \mathcal{J}$ such that
%	\begin{equation*}
%		\| F_{a_i} - F_{a_j} \|_{\infty} \leq \| F_{a_i} - \widetilde{F} \|_{\infty} + \| \widetilde{F}- \widetilde{F}' \|_{\infty} + \| \widetilde{F}' - F_{a_j} \|_{\infty}  \leq 4 \epsilon
%	\end{equation*}
%	for at least one action $a_j \in \mathcal{J}$.
%	% or $I(d) \setminus \mathcal{J}$ is an empty set. 
%	
%	
%	As a result, the procedure described above must terminate in at most $|I(d)| \le n$ steps with $\mathcal{J} \equiv I(d)$.
%	%
%	Moreover, by iteratively employing the triangular inequality, we can conclude that for every pair of actions $a, a' \in I(d)$ it holds $\| F_{a} - F_{a'}\|_{\infty} \le 4 \epsilon n $.
%	%
%	This proves the first of the statement of the lemma.
%	%Finally, since the previous argument holds for each $a_i \in I(d)$, the first part of the proof is concluded.
%	


\boundedactions*
\begin{proof} 
	%
	For ease of presentation, in this proof we adopt the following additional notation.
	%
	Given a set $\mathcal{D}$ of meta-actions and a dictionary $\mathcal{F}$ of empirical distributions computed by means of Algorithm~\ref{alg:action_oracle}, for every meta-action $d \in \mathcal{D}$, we let $\preg_d \subseteq \mathcal{P}$ be the set of all the contracts that have been provided as input to Algorithm~\ref{alg:action_oracle} during an execution in which it computed an empirical distribution that belongs to $\mathcal{F}[d]$.
	%
	Moreover, we let $I(d) \subseteq \mathcal{A}$ be the set of all the agent's actions that have been played as a best response by the agent during at least one of such executions of  Algorithm~\ref{alg:action_oracle}.
	%
	Formally, by exploiting the definition of $\preg_d$, we can write $I(d) \coloneqq \bigcup_{p \in \preg_d} a^\star(p)$.
	%
	%Given a meta-action $d$, we let $\preg_d$ be the set of contracts leading to an observed empirical distribution $\widetilde{F} \in \mathcal{F}[d]$. Additionally, we define $I(d) \subseteq \mathcal{A}$ as the set of associated actions with a meta-action $d$ in a way that $I(d) \coloneqq \cup_{p \in \preg_d} a^\star(p)$.
	
	
	
	First, we prove the following crucial property of Algorithm~\ref{alg:action_oracle}:
	%
	\begin{property}\label{prop_1}
	%
	Given a set $\mathcal{D}$ of meta-actions and a dictionary $\mathcal{F}$ of empirical distributions computed by means of Algorithm~\ref{alg:action_oracle}, under the event $\mathcal{E}_\epsilon$, ti holds that $I(d) \cap I(d') = \varnothing$ for every pair of (different) meta-actions $d \neq  d' \in \mathcal{D}$.
	\end{property}
	%
	In order to show that Property~\ref{prop_1} holds, we assume by contradiction that there exist two different meta-actions $d \neq d' \in \dreg$ such that $a \in I(d)$ and $a \in I(d')$ for some agent's action $a \in \mathcal{A}$, namely that $I(d) \cap I(d') \neq \varnothing$. This implies that there exist two contracts $p \in \preg_{d}$ and $p' \in \preg_{d'}$ such that $a^\star(p) = a^\star(p') = a $. Let $\widetilde{F} \in  \mathcal{F}[d]$ and $\widetilde{F}' \in \mathcal{F}[d']$ be the empirical distributions over outcomes computed by Algorithm~\ref{alg:action_oracle} when given as inputs the contracts $p$ and $p'$, respectively.
	%
	Then, under the event $\mathcal{E}_\epsilon$, we have that $\| \widetilde{F} -\widetilde{F}' \|_\infty \leq 2\epsilon$, since the two empirical distributions are generated by sampling from the same (true) distribution $F_a$.
	%
	Clearly, the way in which Algorithm~\ref{alg:action_oracle} works implies that such empirical distributions are associated with the same meta-action and put in the same entry of the dictionary $\mathcal{F}$.
	%
	This contradicts the assumption that $d \neq d'$, proving that Property~\ref{prop_1} holds. 

	
	
	Next, we show the following additional crucial property of Algorithm~\ref{alg:action_oracle}:
	%
	\begin{property}\label{prop_2}
	%
	Let $\mathcal{D}$ be a set of meta-actions and $\mathcal{F}$ be a dictionary of empirical distributions computed by means of Algorithm~\ref{alg:action_oracle}.
	%
	Suppose that the following holds for every $d \in \dreg$ and $a, a' \in I(d)$:
	\begin{equation}\label{eq:inductive}
		\|F_{a}-F_{a'}\|_{\infty} \le 4 \epsilon (|I(d)|-1).
	\end{equation}
	%
	Then, if Algorithm~\ref{alg:action_oracle} is run again, under the event $\mathcal{E}_\epsilon$ the same condition continues to hold for the set of meta-actions and the dictionary obtained after the execution of Algorithm~\ref{alg:action_oracle}.
	\end{property}
	%
	For ease of presentation, in the following we call $\dreg^{\textnormal{old}}$ and $\mathcal{F}^{\textnormal{old}}$ the set of meta-actions and the dictionary of empirical distributions, respectively, before the last execution of Algorithm~\ref{alg:action_oracle}, while we call $\dreg^{\textnormal{new}}$ and $\mathcal{F}^{\textnormal{new}}$ those obtained after the last run of Algorithm~\ref{alg:action_oracle}.
	%
	Moreover, in order to avoid any confusion, given $\dreg^{\textnormal{old}}$ and $\mathcal{F}^{\textnormal{old}}$, we write $I^{\textnormal{old}}(d)$ in place of $I(d)$ for any meta-action $d \in \dreg^{\textnormal{old}}$.
	%
	Similarly, given $\dreg^{\textnormal{new}}$ and $\mathcal{F}^{\textnormal{new}}$, we write $I^{\textnormal{new}}(d)$ in place of $I(d)$ for any $d \in \dreg^{\textnormal{new}}$.
	
	
	
	%	To prove the lemma we employ an inductive argument showing that if the condition
	%	\begin{equation}%\label{eq:inductive}
	%	 \|F_{a}-F_{a'}\|_{\infty} \le 4 \epsilon (|I(d)|-1)
	%	\end{equation}
	%	holds for all $d$ in $\dreg$ and all $a, a' \in I(d)$ prior to a call to Algorithm~\ref{alg:action_oracle}, then the same property holds after the algorithm's execution. For the ease of notation, we let $\dreg^{\textnormal{new}}$ be the set of meta-actions after a call to Algorithm~\ref{alg:action_oracle}, while we let $\dreg^{\textnormal{old}}$ be the same set before a the call to the same algorithm.
	%	When the set of meta-actions is not updated, \emph{i.e.}, $\dreg^{\textnormal{old}} = \dreg^{\textnormal{new}}$, we use the notation $I^{\textnormal{old}}(d)$ to represent the set of associated actions with a meta-action $d$ before the execution of Algorithm~\ref{alg:action_oracle}. Similarly, $I^{\textnormal{new}}(d)$ represents the same set after the execution of Algorithm~\ref{alg:action_oracle}.
	%
	%	As a preliminary observation, we notice that for each couple of meta-actions $d_i,d_j \in \dreg$ it holds that $I(d_i)\cap I(d_j)=\varnothing$. To show this property, we assume by contradiction that there exist two different meta-actions $d \neq d' \in \dreg$ such that $a \in I(d)$ and $a \in I(d')$ for some agent's action $a \in \mathcal{A}$, meaning that $I(d) \cap I(d') \neq \varnothing$. This implies that there exist two contracts $p \in \preg_{d}$ and $p' \in \preg_{d'}$ such that $a^\star(p) = a^\star(p') = a $. Let $\widetilde{F} \in  \mathcal{F}[d]$ and $\widetilde{F}' \in \mathcal{F}[d']$ be the empirical distributions over outcomes computed by Algorithm~\ref{alg:action_oracle} when given as inputs the contracts $p$ and $p'$, respectively.
	%	%
	%	Under the event $\mathcal{E}_\epsilon$, we have that $\| \widetilde{F} -\widetilde{F}' \|_\infty \leq 2\epsilon$, since the two empirical distributions are generated by sampling from the same (true) distribution $F_a$.
	%	%
	%	Clearly, the way in which Algorithm~\ref{alg:action_oracle} works implies that such empirical distributions are associated with the same meta-action and put in the same entry of the dictionary $\mathcal{F}$.
	%	%
	%	This contradicts the assumption that $d \neq d'$, concluding the proof. 
	%
	% To prove the lemma by induction we observe that the base case is trivially satisfied as the set $\dreg^{\textnormal{new}}$ coincides with a single $\{d\}$ and, thus, there exists a single $I(d)\coloneqq \{a^\star(p)\}$.
	
	
	In order to show that Property~\ref{prop_2} holds, let $p \in \mathcal{P}$ be the contract given as input to Algorithm~\ref{alg:action_oracle} during its last run.
	%
	Next, we prove that, no matter how Algorithm~\ref{alg:action_oracle} updates $\dreg^{\textnormal{old}}$ and $\mathcal{F}^{\textnormal{old}}$ in order to obtain $\dreg^{\textnormal{new}}$ and $\mathcal{F}^{\textnormal{new}}$, the condition in Property~\ref{prop_2} continues to hold.
	%
	We split the proof in three cases.
	%
	%
	% We now consider a generic call to Algorithm~\ref{alg:action_oracle} with $p$ as input and we list all the possible cases in which Algorithm~\ref{alg:action_oracle} updates the dictionary $\mathcal{F}$ and consequently the set $\dreg$.
	%
	\begin{enumerate}
		\item If $|\mathcal{D}^\diamond| = 0$, then $ \dreg^{\textnormal{new}} = \dreg^{\textnormal{old}} \cup \{d^\diamond \}$, where $d^\diamond$ is a new meta-action.
		%
		Since $\mathcal{F}^{\textnormal{new}}$ is built by adding a new entry $\mathcal{F}^{\textnormal{new}}[d^\diamond] = \{ \widetilde{F} \}$ to $\mathcal{F}^{\textnormal{old}}$ while leaving all the other entries unchanged, it holds that $I^{\textnormal{new}}(d^\diamond) = \{ a^\star(p) \}$ and $I^{\textnormal{new}}(d) =  I^{\textnormal{old}}(d)  $ for all $d \in \dreg^{\textnormal{old}}$.
		%
		As a result, the condition in Equation~\eqref{eq:inductive} continues to hold for all the meta-actions $d \in \mathcal{D}^{\textnormal{new}} \setminus \{ d^\diamond\}$.
		%
		Moreover, for the meta-action $d^\diamond$, the following holds:
		\begin{equation*}
			0 = \|F_{a^\star(p)}-F_{a^\star(p)}\|_{\infty} \le 4 \epsilon (|I^{\textnormal{new}} (d^\diamond)|-1) = 0,
		\end{equation*}
		as $|I^{\textnormal{new}} (d^\diamond)|=1$. This proves that the condition in Equation~\eqref{eq:inductive} also holds for $d^\diamond$.
		%
		%We notice that, employing the inductive hypothesis, all the actions $a, a' \in I(d)$ still satisfy the condition $\|F_{a}-F_{a'}\|_{\infty} \le 4 \epsilon (|I(d)|-1)$, for each $d$ belonging to $\dreg^{\textnormal{old}}$.
		%
		%To ensure the validity of Equation~\ref{eq:inductive} after the execution of Algorithm~\ref{alg:action_oracle}, we need to check if such equation is satisfied by the meta-action $d^\diamond$. Noticing that $I(d^\diamond) = \{a^\star(p)\}$, then the following holds:
		%\begin{equation*}
		%	0 = \|F_{a^\star(p)}-F_{a^\star(p)}\|_{\infty} \le 4 \epsilon (|I(d^\diamond)|-1) = 0,
		%\end{equation*}
		%as $|I(d^\diamond)|=1$. This proves that Equation~\ref{eq:inductive} is preserved for all the meta-actions $d$ that belongs to $ \dreg^\textnormal{new}$ after the execution of Algorithm~\ref{alg:action_oracle}.
		
		
		\item If $|\mathcal{D}^\diamond| = 1$, then $ \dreg^{\textnormal{new}} =  \dreg^{\textnormal{old}} $. We distinguish between two cases.
		\begin{enumerate}
			\item In the case in which $a^\star(p) \in \bigcup_{d \in \dreg^{\textnormal{old}}} I^{\textnormal{old}}(d)$, Property~\ref{prop_1} immediately implies that $I^{\textnormal{new}}(d) =  I^{\textnormal{old}}(d)  $ for all $d \in \dreg^{\textnormal{new}}=  \dreg^{\textnormal{old}} $.
			%
			Indeed, if this is \emph{not} the case, then there would be two different meta-actions $d \neq d' \in \dreg^{\textnormal{new}}=  \dreg^{\textnormal{old}} $ such that $I^{\textnormal{new}}(d) =  I^{\textnormal{old}}(d)  \cup \{ a^\star(p) \}$ (since $\mathcal{F}^{\textnormal{new}}[d] = \mathcal{F}^{\textnormal{old}}[d] \cup \{ \widetilde{F} \}$) and $a^\star(p) \in I^{\textnormal{new}}(d') = I^{\textnormal{old}}(d')$, contradicting Property~\ref{prop_1}.
			%
			As a result, the condition in Equation~\eqref{eq:inductive} continues to holds for all the meta-actions after the execution of Algorithm~\ref{alg:action_oracle}.
			%
			% then the set of meta-actions remains unchanged after the execution of Algorithm~\ref{alg:action_oracle}. Moreover, each $I(d)$ remains unchanged as well. Consequently, Equation~\ref{eq:inductive} is still verified after the execution of Algorithm~\ref{alg:action_oracle}.
			
			
			\item
			%If $a^\star(p) \notin \cup_{d \in \dreg^{\textnormal{old}}} I(d)$, under the event $\mathcal{E}_\epsilon$, we have $\| F_{a^\star(p) }- \widetilde{F} \|_{\infty} \le \epsilon$ and there exists a $F \in \mathcal{F}[d]$ for some $d$ such that $\| F - \widetilde{F} \|_{\infty} \le 2\epsilon$. This is ensured by observing how Algorithm~\ref{alg:action_oracle} works. Furthermore, there exists an action $a \in I(d)$ such that $\| F - F_a  \|_{\infty} \le \epsilon$ as long as the event $\mathcal{E}_\epsilon $ holds. Then by employing a triangular inequality we get:
			In the case in which $a^\star(p) \notin \bigcup_{d \in \dreg^{\textnormal{old}}} I^{\textnormal{old}}(d)$, the proof is more involved. 
			%
			Let $d \in \dreg^{\textnormal{new}} =  \dreg^{\textnormal{old}}$ be the (unique) meta-action in the set $\dreg^\diamond$ computed by Algorithm~\ref{alg:action_oracle}.
			%
			Notice that $I^{\textnormal{new}}(d) = I^{\textnormal{old}}(d) \cup \{a^\star(p)\}$ by how Algorithm~\ref{alg:action_oracle} works.
			%
			As a first step, we show that, for every pair of actions $a,a' \in I^{\textnormal{new}}(d)$ with $a = a^\star(p)$, Equation~\eqref{eq:inductive} holds.
			%
			Under the event $\mathcal{E}_\epsilon$, it holds that $\| F_{a^\star(p) }- \widetilde{F} \|_{\infty} \le \epsilon$, where $\widetilde{F}$ is the empirical distribution computed by Algorithm~\ref{alg:action_oracle}.
			%
			Moreover, since the meta-action $d$ has been added to $\dreg^\diamond$ by Algorithm~\ref{alg:action_oracle}, there exists an empirical distribution $F \in \mathcal{F}[d]$ such that $\| \widetilde{F} - {F} \|_{\infty} \le 2\epsilon$, and, under the event $\mathcal{E}_\epsilon$, there exists an action $a'' \in I^\textnormal{old}(d)$ such that $\| F - F_{a''} \|_{\infty} \le \epsilon$.
			%
			Then, by applying the triangular inequality we can show that:
			\begin{equation*}
				\| F_{a^\star(p) }- F_{a''}\|_{\infty} \le \| F_{a^\star(p) }- \widetilde{F} \|_{\infty}  + \|  \widetilde{F} - F\|_{\infty} + \| F - F_{a''}  \|_{\infty} \le 4 \epsilon.
			\end{equation*}
			By using the fact that the condition in Equation~\eqref{eq:inductive} holds for $\dreg^{\textnormal{new}}$ and $\mathcal{F}^{\textnormal{new}}$, for every action $a' \in I^{\textnormal{new}}(d)$ it holds that:
			\begin{align*}
				\| F_{a^\star(p)}- F_{ a'} \|_{\infty} & \le \| F_{a^\star(p)} - F_{a''}\|_{\infty} + \| F_{a''}- F_{a'}  \|_{\infty} \\
				&  \le 4 \epsilon (|I^{\textnormal{old}}({d}) |-1 ) + 4 \epsilon  \\ &  \le 4 \epsilon | I^{\textnormal{old}}(d)| \\
				&= 4 \epsilon (|I^{\textnormal{new}}(d) |  - 1),
			\end{align*}
			where the last equality holds since $|I^{\textnormal{new}}(d) | = |I^{\textnormal{old}}(d) | +1$, as $a^\star(p) \notin \bigcup_{d \in \dreg^{\textnormal{old}}} I^{\textnormal{old}}(d)$.
			%
			This proves that Equation~\eqref{eq:inductive} holds for every pair of actions $a,a' \in I^{\textnormal{new}}(d)$ with $a = a^\star(p)$.
			%
			For all the other pairs of actions $a,a' \in I^{\textnormal{new}}(d)$, the equation holds since it was already satisfied before the last execution of Algorithm~\ref{alg:action_oracle}.
			%
			Analogously, given that $I^{\textnormal{new}}(d) = I^{\textnormal{old}}(d)$ for all the meta-actions in $\dreg^{\textnormal{new}} \setminus \{d\}$, we can conclude that the condition in Equation~\eqref{eq:inductive} continues to holds for such meta-actions as well.
			%&
			%
			% Consequently for each couple of actions $a',a'' \in I^{\textnormal{new}}(d)$ we have that: $ \| F_{a' }- F_{a'' }\|_{\infty} \le 4 \epsilon | I^{\textnormal{old}}(d)| \le 4 \epsilon (| I^{\textnormal{new}}(d)|-1)$.  Finally, observing that all the remaining sets $I(d')$ with $d'$ in $\dreg^{\textnormal{new}} \setminus \{d\}$ remain unchanged, we have that Equation~\ref{eq:inductive} is still satisfied after the execution of Algorithm~\ref{alg:action_oracle} for all the meta-actions belonging to $\dreg^{\textnormal{new}}$ .
		\end{enumerate}
		
		\item  If $|\mathcal{D}^\diamond| > 1$, then $\mathcal{D}^{\textnormal{new}} = \left( \mathcal{D}^{\textnormal{old} } \setminus \mathcal{D}^\diamond \right) \cup \{ d^\diamond \}$, where $d^\diamond$ is a new meta-action.
		%
		Clearly, the condition in Equation~\eqref{eq:inductive} continues to hold for all the meta-actions in $ \mathcal{D}^{\textnormal{old} } \setminus \mathcal{D}^\diamond  $.
		%
		We only need to show that the condition holds for $d^\diamond$.
		%
		We distinguish between two cases.
		%
		\begin{enumerate}
			
			\item In the case in which $a^\star(p) \in  \bigcup_{d \in \dreg^{\textnormal{old}}} I^{\textnormal{old}}(d)$, let us first notice that $a^\star(p) \in I^{\textnormal{old}}(d)$ for some $d \in \dreg^\diamond$, otherwise Property~\ref{prop_1} would be violated (as $a^\star(p) \in I^{\textnormal{new}}(d^\diamond)$ by definition).
			%
			Moreover, it is easy to see that $I^\textnormal{new}(d^\diamond)= \bigcup_{d \in \dreg^{\diamond} } I^\textnormal{old}(d)$ and, additionally, $I^\textnormal{old}(d) \cap I^\textnormal{old}(d') = \varnothing$ for all $d \neq d' \in \dreg^\diamond$, given how Algorithm~\ref{alg:action_oracle} works and thanks to Property~\ref{prop_1}.
			%
			In the following, for ease of presentation, we assume w.l.o.g.~that the meta-actions in $\dreg^\diamond$ are indexed by natural numbers so that $\dreg^{\diamond} \coloneqq \{ d_1, \dots d_{|\mathcal{D}^\diamond|}\}$ and $a^\star(p) \in I^{\textnormal{old}}(d_1)$.
			%
			Next, we show that Equation~\eqref{eq:inductive} holds for every pair $a, a' \in I^\textnormal{new}(d^\diamond)$.
			%
			% we assume w.l.o.g.~that $a^\star(p) \in I^{\textnormal{old}}(d_1)$ for some $d_1 \in \dreg^{\textnormal{old}}$ and $\dreg^{\diamond} \coloneqq \{ d_1, \dots d_{|\mathcal{D}^\diamond|}\}$ and we notice that $I^\textnormal{new}(d^\diamond)= \cup_{d \in \dreg^{\diamond} } I^\textnormal{old}(d)$.
			%
			First, by employing an argument similar to the one used to prove Point~2.2, we can show that, for every $d_j \in \dreg^\diamond$ with $j > 1$, there exists an action $a'' \in I^{\textnormal{old}}(d_j)$ such that $\| F_{a^\star(p) }- F_{a''}\|_{\infty} \le 4 \epsilon$. Then, for every pair of actions $a, a' \in I^\textnormal{new}(d^\diamond)$ such that $a \in I^{\textnormal{old}}(d_1)$ and $a' \in I^{\textnormal{old}}(d_j)$ for some $d_j \in \dreg^\diamond$ with $j > 1$, the following holds:
			\begin{align*}
				\| F_{a }- F_{a'}\|_{\infty} & \le \| F_{a}  - F_{a^\star(p) } \|_{\infty} + \| F_{a^\star(p) }- F_{a''}  \|_{\infty} + \|  F_{a''} - F_{a'}  \|_{\infty}\\
				&   \le 4 \epsilon( | I^{\textnormal{old}}(d_1)| -1 ) + 4 \epsilon + 4 \epsilon( | I^{\textnormal{old}}({d_j}) | -1 )\\
				&   = 4 \epsilon( | I^{\textnormal{old}}({d_1})| + | I^{\textnormal{old}}({d_j} ) |-1 )  \\
				&  \leq 4 \epsilon( | I^{\textnormal{new}}(d^\diamond) | -1 ),	 	
			\end{align*}
			where the first inequality follows from an application of the triangular inequality, the second one from the fact that Equation~\eqref{eq:inductive} holds before the last execution of Algorithm~\ref{alg:action_oracle}, while the last inequality holds since $| I^{\textnormal{new}}(d^\diamond) |= \sum_{d \in \dreg^\diamond}| I^{\textnormal{old}}(d) |$ given that $I^\textnormal{old}(d) \cap I^\textnormal{old}(d') = \varnothing$ for all $d \neq d' \in \dreg^\diamond$ thanks to Property~\ref{prop_1}. 
			%
			% follows by employing the inductive hypothesis and the second inequality holds by observing that $I_{d^\diamond} = \sum_{d \in \dreg^\diamond} I_d$ as $I(d_i) \cap I(d_j)=\varnothing$. 
			%
			%
			As a final step, we show that Equation~\eqref{eq:inductive} holds for every pair of actions $a, a' \in I^\textnormal{new}(d^\diamond)$ such that $a \in I^{\textnormal{old}}(d_i)$ and $a' \in I^{\textnormal{old}}(d_j)$ for some $d_i, d_j \in \dreg^\diamond$ with $i \neq j > 1$.
			%
			By the fact that $d_i, d_j \in \dreg^\diamond$ and triangular inequality, it follows that there exist $a'' \in I^{\textnormal{old}}(d_i)$ and $a''' \in I^{\textnormal{old}}(d_j)$ such that $\|F_{a''} -\widetilde F \|_{\infty} \le 3 \epsilon$ and $\| F_{a'''} - \widetilde F \|_{\infty} \le 3 \epsilon$ under $\mathcal{E}_\epsilon$.
			%
			Then, with steps similar to those undertaken above, we can prove the following:
			%
			%As a final step we need to show that for each $a' \in I^{\textnormal{old}}(d_i)$ and $a'' \in I^{\textnormal{old}}(d_j)$ with $i,j \in \{2, \dots, |\mathcal{D}^\diamond| \}$ it holds $\|F_{a'}-F_{a''}\|_{\infty} \le 4 \epsilon (|I(d^\diamond)|-1)$. By employing the triangular inequality we can show that there exists $a_i $ that belongs to $I(d_i)$ and $a_j$ that belongs to $I(d_j)$ such that $\|\widetilde F - F_{a_i}\|_{\infty} \le 3 \epsilon$ and $\| \widetilde F - F_{a_j}\|_{\infty} \le 3 \epsilon$. Then, for each actions $a'$ that belongs to $ I(d_i)$ and action $a'' $ that belongs to  $I(d_j)$, we have:
			%
			\begin{align*}
				\| F_{a}- F_{a'}\|_{\infty} & \le\| F_{a}- F_{a''}  \|_{\infty} +\| F_{a''} -\widetilde F \|_{\infty} + \| \widetilde F - F_{a'''}  \|_{\infty}  + \|  F_{a'''} - F_{a'}  \|_{\infty}\\
				&   \le 4 \epsilon(|I^{\textnormal{old}}({d_i}) | -1 ) + 3\epsilon +  3 \epsilon + 4 \epsilon(|I^{\textnormal{old}}({d_j}) | -1 )\\
				&   = 4 \epsilon(| I^{\textnormal{old}}({d_i})| +  |I^{\textnormal{old}}({d_j}) | -1 ) +2\epsilon  \\
				&\le 4 \epsilon(|I^{\textnormal{new}}({d^\diamond})| -1 ),
			\end{align*}
			which shows that Equation~\ref{eq:inductive} is satisfied for all the pairs of actions $a,a' \in I^\textnormal{new}(d^\diamond)$.
			%
			%Showing that Equation~\ref{eq:inductive} is satisfied for all the actions $a,a' \in I^\textnormal{new}(d^\diamond) $.
			
		
			\item In the case in which $a^\star(p) \notin \bigcup_{d \in \dreg^{\textnormal{old}}} I^{\textnormal{old}}(d)$, let us first observe that $I^{\textnormal{new}}(d^\diamond) =  \bigcup_{d \in\dreg^\diamond } I^{\textnormal{old}}(d) \cup \{a^\star(p)\}$ and $I^\textnormal{old}(d) \cap I^\textnormal{old}(d') = \varnothing$ for all $d \neq d' \in \dreg^\diamond$, given how Algorithm~\ref{alg:action_oracle} works and thanks to Property~\ref{prop_1}.
			%
			Next, we show that Equation~\ref{eq:inductive} holds for every pair of actions $a,a' \in  I^\textnormal{new}(d^\diamond)$.
			%
			As a first step, we consider the case in which $a = a^\star(p)$, and we show that Equation~\ref{eq:inductive} holds for every $a' \in I^\textnormal{new}(d^\diamond)$ such that $a' \in I^\textnormal{old}(d)$ for some $d \in \mathcal{D}^\diamond$.
			%
			In order to show this, we first observe that, given how Algorithm~\ref{alg:action_oracle} works, there exists $F \in \mathcal{F}^\textnormal{old}[d]$ such that $ \| \widetilde F - F \| \le 2 \epsilon $, and, under the event $\mathcal{E}_\epsilon$, there exists an action $a'' \in I^\textnormal{old}(d)$ such that $ \|  F_{a''} - F \| \le \epsilon $.
			%
			Then:
			%
			%
			%then $I^{\textnormal{old}}(d^\diamond)\coloneqq \cup_{d \in\ \dreg^\diamond } I(d) \cup \{a^\star(p)\}$. Consequently  we need to show that for each $a,a' $ that belongs to $ I(d^\diamond)$ we have $\|F_{a}-F_{a'}\|_{\infty} \le 4 \epsilon (|I(d^\diamond)|-1)$ since the sets $I(d)$ with $d$ belonging to $\dreg^\textnormal{old} \setminus \dreg^\diamond$ remains unchanged.
			%
			%To do that we fist show that for each $a' \in I(d)$ and $d \in \dreg^\diamond$ it holds $\|F_{a^\star(p)}-F_{a'}\|_{\infty} \le 4 \epsilon (|I(d^\diamond)|-1)$. To prove that we observe that there exists a $F_d $ that belongs to $\mathcal{F}[d]$ such that $ \| \widetilde F - F_d \| \le 2 \epsilon $ as guaranteed by Algorithm~\ref{alg:action_oracle}, and for each eta action $d \in \dreg^\diamond$ there exists an action $a \in I(d)$ such that $ \|  F_a - F_d \| \le \epsilon $, under the event $\mathcal{E}_\epsilon$. Then by employing the triangular inequality and assuming the event $\mathcal{E}_\epsilon$ holds, for each $a' \in I(d)$ the following inequality holds:
			%
			\begin{align*}
				\| F_{a^\star(p) } - F_{a'}\|_{\infty} &\le\| F_{a^\star(p) } - \widetilde F  \|_{\infty} +\| \widetilde F - F\|_{\infty} + \| F- F_{a''}  \|_{\infty}  + \| F_{a''}- F_{a'}  \|_{\infty} \\
				&     \le 4 \epsilon +  4 \epsilon(|I^{\textnormal{old}}({d}) | -1 ) \\
				&   \le 4 \epsilon(|I^{\textnormal{new}}(d^\diamond)| -1 ),
			\end{align*}
			where the last inequality holds since $|I^{\textnormal{new}}_{d^\diamond}| = \sum_{d \in \dreg^\diamond} |I^{\textnormal{old}}(d)|$ as $I^{\textnormal{old}}(d) \cap I^{\textnormal{old}}(d')=\varnothing$ for all $d \neq d' \in \dreg^\diamond$. 
			%
			%
			% the first inequality above follows by employing the inductive hypothesis and the second inequality holds by observing that $I_{d^\diamond} = \sum_{d \in \dreg^\diamond} I_d$ as $I(d_i) \cap I(d_j)=\varnothing$. 
			%
			%
			As a second step, we consider the case in which $a \in I^{\textnormal{old}}(d)$ and $a' \in I^{\textnormal{old}}(d')$ for some pair $d \neq d' \in \dreg^\diamond$.
			%
			Given how Algorithm~\ref{alg:action_oracle} works, there exist $F \in \mathcal{F}^{\textnormal{old}}[d]$ and $F' \in \mathcal{F}^{\textnormal{old}}[d']$ such that $\| F -\widetilde F \| \le 2 \epsilon$ and $\| F' -\widetilde F \| \le 2 \epsilon$.
			%
			Moreover, by the fact that $d, d' \in \dreg^\diamond$ and the triangular inequality, it follows that there exist $a'' \in I^{\textnormal{old}}(d)$ and $a''' \in I^{\textnormal{old}}(d')$ such that $\|F_{a''} -\widetilde F \|_{\infty} \le 3 \epsilon$ and $\| F_{a'''} - \widetilde F \|_{\infty} \le 3 \epsilon$ under the event $\mathcal{E}_\epsilon$.
			%
			Then, the following holds:
			%
			%As a second step, we show that for each action $a' \in I(d_i)$ and $a'' \in I(d_j)$ we have $\| F_{a'} - F_{a''}\| \le 4 \epsilon (|I(d^\diamond)|-1)$.
			%To show that we know that there exist  $F_i$ that belongs to $\mathcal{F}[d_i]$ and $F_j$ that belongs to $\mathcal{F}[d_j]$ such that $\| F_j -\widetilde F \| \le 2 \epsilon$ and $\| F_i -\widetilde F \| \le 2 \epsilon$. Then, for each actions $a'$ that belongs to $ I(d_i)$ and action $a'' $ that belongs to  $I(d_j)$, we have:
			%
			\begin{align*}
				\| F_{a}- F_{a'}\|_{\infty} & \le\| F_{a}- F_{a''}  \|_{\infty} +\| F_{a''} -\widetilde F \|_{\infty} + \| \widetilde F - F_{a'''}  \|_{\infty}  + \|  F_{a'''} - F_{a'}  \|_{\infty}\\
				&   \le 4 \epsilon(|I^{\textnormal{old}}({d}) | -1 ) + 3\epsilon +  3 \epsilon + 4 \epsilon(|I^{\textnormal{old}}({d'}) | -1 )\\
				&   = 4 \epsilon(|I^{\textnormal{old}}({d}) | +  | I^{\textnormal{old}}(d') | -1 ) +2\epsilon \\
				& \le 4 \epsilon(|I^{\textnormal{new}}({d^\diamond}) | -1 ).
			\end{align*}
		\end{enumerate}
	\end{enumerate}


	Finally, by observing that the size of $I(d)$ for each possible set of meta-actions $\dreg$ is bounded by the number of agent's actions $n$, for each $a,a' \in I(d)$ and $d \in \dreg$ we have that:
	\begin{equation}\label{eq:rec_finale}
		\lVert F_{a}- F_{a'} \rVert\le 4n\epsilon.
	\end{equation}
	Then, under the event $\mathcal{E}_\epsilon$, by letting $a' \in I(d)$ be the action leading to the empirical distribution $\widetilde F_d$ computed at Line~\ref{line:assoc_ditrib} in Algorithm~\ref{alg:action_oracle} we have $ \| \widetilde F_d -F_{ a' } \|_\infty\le \epsilon$.
	Finally, combining the two inequalities, we get:
	\begin{equation*}
		\|\widetilde F_{d}- F_{a^\star(p)}\|_{\infty} \leq \|\widetilde F_d- F_{ a' } \|_{\infty} + \|F_{ a ' }- F_{a^\star(p)} \|_{\infty}  \le 4\epsilon n+ \epsilon\le 5\epsilon n.
	\end{equation*}
	This implies that $a^\star(p) \in A(d)$.
	
	To conclude the proof we show that $I(d)\subseteq A(d)$. Let $a' \in I(d)$ be an arbitrary action, and let $\widetilde F_d$ be the empirical distribution computed by Algorithm~\ref{alg:action_oracle} by sampling from $F_a$. Then for each $a'' \in I(d)$ we have:
	\begin{equation*}
		\|\widetilde F_{d}- F_{a''}\|_{\infty} \leq \|\widetilde F_d- F_{ a' } \|_{\infty} + \|F_{ a ' }- F_{a''} \|_{\infty}  \le 4\epsilon n+ \epsilon\le 5\epsilon n.
	\end{equation*}
	Since the latter argument holds for all the actions $a \in I(d)$ this hows that 
		\begin{equation}
		\lVert F_{a}- F_{a'} \rVert_{\infty} \le 4n\epsilon,
	\end{equation}
	for all the actions $a,a' \in A(d)$.
	\end{proof}

\finaliterations*
\begin{proof}
	
	As a first step, we observe that under the event $\mathcal{E}_\epsilon$, the size of $\dreg$ increases whenever the principal observes an empirical distribution $\widetilde{F}$ that satisfies $\| \widetilde{F} - \widetilde{F}' \|_{\infty} \ge 2 \epsilon$ for every $\widetilde{F}' \in \mathcal{F}$. This condition holds when the agent chooses an action that has not been selected before. This is because, if the principal commits to a contract implementing an action the agent has already played, the resulting estimated empirical distribution $\widetilde{F}$ must satisfy $\| \widetilde{F} - F \|_{\infty} \le 2 \epsilon$ for some $d \in \mathcal{D}$ and $F \in \mathcal{F}_{d}$, as guaranteed by the event $\mathcal{E}_{\epsilon}$. Consequently, the cardinality of $\mathcal{D}$ can increase at most $n$ times, which corresponds to the total number of actions. Furthermore, we observe that the cardinality of $\dreg$ can decrease by merging one or more meta-actions into a single one, with the condition $|\dreg| \geq 1$. Therefore, in the worst case the cardinality of $\mathcal{D}$ is reduced by one $n$ times, resulting in $2 n $ being the maximum number of times Algorithm~\ref{alg:action_oracle} returns $\bot$.
\end{proof}



\costactions*
\begin{proof}
	%Fix a meta-action $d \in \mathcal{D}$ and assume that the event $\mathcal{E}_\epsilon$ holds. Let $a' \in \argmin_{a \in A(d)} c_a $ while let $p'$ be a contract such that $a^\star(p')=a'$. Furthermore let $a \in A(d)$ then there exists a $p \in \preg$ such that $a^\star(p)=a$. Then we have:
	Let $d \in \mathcal{D}$ be a meta-action and assume that the event $\mathcal{E}_\epsilon$ holds. Let $a' \in \argmin_{a \in A(d)} c_a$ and let $p'$ be a contract such that $a^\star(p')=a'$. By employing Definition~\ref{def:asssoc_actions}, we know that for any action $a \in A(d)$, there exists a contract $p \in \preg$ such that $a^\star(p)=a$. Then the following inequalities hold:
	\begin{equation}\label{eq:cost_1}
		4B \epsilon m  n \ge \| {F}_{a} - {F}_{a'} \|_{\infty} \| p\|_{1}  \ge \sum_{\omega \in \Omega} \left( {F}_{a,\omega} - {F}_{a',\omega} \right) p_{\omega}  \ge  {c}_{a} - {c}_{a'},
	\end{equation} 
	and, similarly, 
	\begin{equation}\label{eq:cost_2}
		4B \epsilon m  n \ge \| {F}_{a'} - {F}_{a} \|_{\infty} \| p'\|_{1}  \ge \sum_{\omega \in \Omega} \left( {F}_{a',\omega} - {F}_{a,\omega} \right) p'_{\omega}  \ge  {c}_{a'} - {c}_{a} . 
	\end{equation} 
	In particular, in the chains of inequalities above, the first `$\geq $' holds since $\| F_{a} - {F}_{a'}\|_{\infty} \le 4 \epsilon n$ by Lemma~\ref{lem:boundedactions} and the fact that the norm $\|\cdot\|_1$ of contracts in $\mathcal{P}$ is upper bounded by $Bm$.
	% observing that $\|p\|_1 \le Bm$.
	%
	Finally, by applying Equations~\eqref{eq:cost_1}~and~\eqref{eq:cost_2}, we have that:
	\begin{equation*}
	|  {c}_{a} - {c}_{a'} | =	|  {c}_{a} - c(d) |   \le 4 B \epsilon m  n, 
	\end{equation*} 
	for every possible action $a \in A(d)$ since $c_{a'}=c(d)$ by definition, concluding the proof. 
\end{proof}


\closeutility*
\begin{proof}
To prove the lemma we observe that for each action $a \in A(d)$ it holds:
\begin{align*}
	\left|\sum_{\omega \in \Omega} \widetilde{F}_{d,\omega} p_\omega -c(d) - \sum_{\omega \in \Omega} F_{a, \omega} p_\omega + c_{a} \right| & \le \| \widetilde{F}_{d } - {F}_{a} \|_{\infty} \| p \|_{1}   +  |c(d)-c_a|\\
	& \le 5 B \epsilon m  n  +  4 B \epsilon m  n = 9B \epsilon m  n,
\end{align*}
%where the first inequality holds employing the triangular and the holder's inequality, while the second inequality holds by means of Lemma and employing the Definition ensuring that  $\| \widetilde{F}_{d } - {F}_{a} \|_{\infty} \le 5 \epsilon n$, for each action $a \in A(d)$ and observing that the norm $\|\cdot\|_1$ of contracts in $\mathcal{P}$ is upper bounded by $Bm$.

where the first inequality holds by applying the triangular inequality and Holder's inequality. The second inequality holds by employing Lemma~\ref{lem:cost_actions} and leveraging Definition~\ref{def:asssoc_actions} that guarantees $\| \widetilde{F}_{d} - {F}_{a} \|_{\infty} \le 5 \epsilon n$ for every action $a \in A(d)$. Additionally, we observe that the $\|\cdot\|_1$ norm of contracts in $\mathcal{P}$ is upper bounded by $Bm$, which concludes the proof.
\end{proof}


\begin{comment}
%\actintersection*
\begin{proof}
	%	When either $\mathcal{D}^\diamond = \varnothing$ or $|\mathcal{D}^\diamond| > 1$, the property is clearly preserved by Algorithm~\ref{alg:action_oracle} if it held before the algorithm execution.
	%	%
	%	The only case in which Algorithm~\ref{alg:action_oracle} may \emph{not} preserve the property is when $|\mathcal{D}^\diamond| = 1$.
	%
	By contradiction, suppose that there exist two different meta-actions $d \neq d' \in \dreg$ such that $a \in I(d)$ and $a \in I(d')$ for some agent's action $a \in \mathcal{A}$, which means that $I(d) \cap I(d') \neq \varnothing$. By definition of associated action (Definition~\ref{def:associated_actions}), this implies that there exist two contracts $p \in \preg_{d}$ and $p' \in \preg_{d'}$ such that $a^\star(p) = a^\star(p') = a $.
	%
	Let $\widetilde{F} \in \Delta_{\Omega}$ and $\widetilde{F}' \in \Delta_{\Omega}$ be the empirical distributions over outcomes computed by Algorithm~\ref{alg:action_oracle} when given as inputs the contracts $p$ and $p'$, respectively.
	%
	Under the event $\mathcal{E}_\epsilon$, we have that $\| \widetilde{F} -\widetilde{F}' \|_\infty \leq 2\epsilon$, since the two empirical distributions are generated by sampling from the same (true) distribution $F_a$.
	%
	Clearly, the way in which Algorithm~\ref{alg:action_oracle} works implies that such empirical distributions are associated with the same meta-action and put in the same entry of the dictionary $\mathcal{F}$.
	%
	This contradicts the assumption that $d \neq d'$, concluding the proof. 
	%	
	%	$p$ and $p'$ belong to $\mathcal{F}_d$ for some $d \in \dreg$ by the construction of the \texttt{Action-Oracle} algorithm. However, this leads to a contradiction with the fact that $p^1 \in \preg_{d_1}$ and $p^2 \in \preg_{d_2}$, with $d_1 \neq d_2$, concluding the proof.
\end{proof}
\end{comment}

%We conclude the subsection by providing some auxiliary lemmas related to the \texttt{Action-Oracle} procedure that will be useful  in the following.
%
%The first lemma will be employed to bound the probability of the clean event $\mathcal{E}_\epsilon$ in the proof of Theorem~\ref{thm:finalthm}.
%
%The second lemma will be useful in proving Lemma~\ref{lem:epsilonbr} in the following subsection.


We conclude the subsection by providing an auxiliary lemmas related to the \texttt{Action-Oracle} procedure that will be employed to bound the probability of the clean event $\mathcal{E}_\epsilon$ in the proof of Theorem~\ref{thm:finalthm}.

\begin{lemma}\label{lem:hoeffding}
	Given $\alpha, \epsilon \in (0,1)$, let $\widetilde{F} \in \Delta_{\Omega}$ be the empirical distribution computed by Algorithm~\ref{alg:action_oracle} with $q:=\left\lceil \frac{1}{2 \epsilon^2} \log \left(\frac{2 m}{\alpha}\right)\right\rceil $ for a contract $p \in \mathcal{P}$ given as input.
	%
	Then, it holds that $\mathbb{E}[\widetilde{F}_\omega]=F_{a^\star(p),\omega} $ for all $\omega \in \Omega $ and, with probability of at least $1-\alpha$, it also holds that $ \|\widetilde{F}-F_{a^\star(p)} \|_{\infty} \le \epsilon $.
	%	
	%	Given $\alpha, \epsilon >0$ and a probability distribution $F \in \Delta_{\Omega}$, we let $\widetilde{F} \in \Delta_{\Omega}$ be the empirical distribution computed by Algorithm \ref{alg:action_oracle} with $q=\left\lceil \frac{1}{2 \epsilon^2} \log \left(\frac{2 m}{\alpha}\right)\right\rceil $ rounds. Then we have $\mathbb{E}[\widetilde{F}_\omega]=F_\omega $ for each $\omega \in \Omega $ and with probability of at least $1-\alpha$ it holds $ \|\widetilde{F}-F \|_{\infty} \le \epsilon $.
\end{lemma}
\begin{proof}
	By construction, the empirical distribution $\widetilde{F} \in \Delta_\Omega$ is computed from $q$ i.i.d.~samples drawn according to $F_{a^\star(p)}$, with each $\omega \in \Omega$ being drawn with probability $F_{a^\star(p), \omega}$.
	%
	Therefore, the empirical distribution $\widetilde{F}$ is a random vector supported on $\Delta_\Omega$, 
	whose expectation is such that $\mathbb{E}[\widetilde{F}_\omega]=F_{a^\star(p),\omega}$ for all $\omega \in \Omega$.
	%
	Moreover, by Hoeffding’s inequality, for every $\omega \in \Omega$, we have that:
	\begin{equation}\label{eq:hoeffding}
		\mathbb{P}\left\{ |\widetilde{F}_\omega-\mathbb{E}[\widetilde{F}_\omega]| \ge \epsilon \right\} =\mathbb{P} \left\{ |\widetilde{F}_\omega-F_{a^\star(p), \omega}| \ge \epsilon \right\} \ge 1- 2e^{-2q\epsilon^2}.
	\end{equation}
	Then, by employing a union bound and Equation~\eqref{eq:hoeffding} we have that:
	\begin{equation*}
		\mathbb{P}\left\{  \| \widetilde{F} - F_{a^\star(p)} \|_{\infty} \le \epsilon \right\} = \mathbb{P} \left\{  \bigcap_{\omega \in \Omega } \left\{ |\widetilde{F}_\omega-F_{a^\star(p), \omega}| \le \epsilon \right\} \right\} \ge 1 - 2 m e^{-2q\epsilon^2} \ge 1-\alpha,
	\end{equation*}
	where the last inequality holds by definition of $q$.
\end{proof}




\begin{comment}
		\begin{proof} 
		In the following, we denote with $\widetilde{F}_{a_i,a_j} \in \mathcal{F}_{d}$ the empirical distribution observed by the principal, satisfying $\| \widetilde{F}_{a_i,a_j} - \widetilde{F}_{i} \|_{\infty} \le 2\epsilon$ and $\| \widetilde{F}_{a_i,a_j} - \widetilde{F}_{j}\|_{\infty} \le 2\epsilon$ with $\widetilde{F}_i, \widetilde{F}_j \in \mathcal{F}_{d}$ empirical distributions observed by the principal, where $\widetilde{F}_{i} \in B_\epsilon(F_{a_i})$ and $\widetilde{F}_{j} \in B_\epsilon(F_{a_{j}})$. We define the set $B_\epsilon(F)$ as follows:
		\begin{equation*}
			B_\epsilon(F) \coloneqq \{F' \in \Delta_\Omega \, | \, \| F - F' \|_{\infty} \le \epsilon   \}.
		\end{equation*}
		Notice that, under the event $\mathcal{E}_{\epsilon}$, we know that $\widetilde{F}_i \in B_\epsilon(F_{a_i})$ and $\widetilde{F}_{j} \in B_\epsilon(F_{a_{j}})$.
		%\alb{Sopra non si capisce cosa sono quelle distribuzioni.}
		We observe that, for each couple of actions $a_{i_1}, a_{i_2} \in I(d)$ there exists a $I'=\{a_{i_1}, \dots, a_{i_2} \} \subseteq I(d)$ and a sequence of empirical distributions $\{ \widetilde{F}_{a_i,a_{i+1}}\}_{i \in I'} \subseteq \mathcal{F}_{d}$ define as above. This is because \texttt{UPDATE-ACTIONS} identifies action $a_{i_1}$ and action $ a_{i_2}$ as a single $d \in \mathcal{D}$. We first we observe that for each $a_i \in I'$ it holds:
		\begin{align*}
			\| F_{a_{i}} - F_{a_{i+1}}\|_{\infty} & =   \| F_{a_{i}} \pm \widetilde{F}_{a_i,a_{i+1}} \pm \widetilde{F}_{i+1} +  F_{a_{i+1}}\|_{\infty} \\
			& =   \| F_{a_{i}} - \widetilde{F}_{a_i,a_j} \|_{\infty}  + \| \widetilde{F}_{a_i,a_j} - \widetilde{F}_{i+1}\|_{\infty} + \| \widetilde{F}_{i+1} - F_{a_{i+1}} \|_{\infty} \\
			& \le 4 \epsilon,
		\end{align*}
		where the inequality holds under the event $\mathcal{E}_\epsilon$ and assuming without loss of generality that $\widetilde{F}_{a_i,a_j} \in B_\epsilon(F_{a_i})$. Then we have:
		\begin{equation*}
			\| F_{a_{i_1}} - F_{a_{i_2}} \|_{\infty} = \| \sum_{i \in I'} F_{a_{i}} - F_{a_{i+1}} \|_{\infty} \le \sum_{i \in I'} \| F_{a_{i}} - F_{a_{i+1}}\|_{\infty} \le 5 \epsilon |I'| \le 4 \epsilon n,
		\end{equation*}
		where we let $F_{a_{i_{2}+1}}= F_{a_{i_2}}$ for the ease of notation. Finally, we observe that for each $\widetilde{F} \in \mathcal{F}_d$ we have for each action $a_i \in I(d)$:
		\begin{equation*}
			\| F_{a_{i}} - \widetilde{F} \|_{\infty} \le  \|  F_{a_{i}}  -  F_{a_{j}}  \|_{\infty} +  \|  \widetilde{F}   -  F_{a_{j}}  \|_{\infty}  \le 5 \epsilon n,
		\end{equation*}
		where we let $\widetilde{F} \in B_\epsilon(F_{a_j})$, concluding the proof.
	\end{proof}
	%To prove the lemma we  a constructive argument which is the rationale behind the \texttt{Action-Oracle} algorithm. We start by picking at random an action $a_i \in I(d)$ and we initialize a set $\mathcal{J} \coloneqq \{ a_i \}$. Then we prove that either $\mathcal{J} = I(d)$ or there exists a non emtpy subset $I \subseteq I(d) \setminus \mathcal{J}$ and a $F \in F_d $ such that $\| F - F_{a_j}\|_{\infty} \le 2 \epsilon$ and $\| F - F_{a_i}\|_\infty \le 2 \epsilon$. Finally we update the set $\mathcal{J} = \mathcal{J} \cup I$ and we iterate the procedure described above until $\mathcal{J}=I(d)$. .. We summarize the procedure in Algorithm 1.
\end{comment}



%
